

We now construct a trace whose six beam strain measures are \emph{energetically
conjugate} to the classical Saint-Venant resultants. 
Throughout this section we
work entirely at the Saint-Venant level and use only explicitly defined
quantities.

\paragraph{Saint-Venant pointwise strains.}

For a prismatic bar with axial unit vector $\FrameBasis$ and cross-section
$\Omega \subset \mathbb{R}^3$ (lying in the plane orthogonal to $\FrameBasis$), the
Saint-Venant solution is parameterized by
\[
\bm{p} =
\begin{pmatrix}
a_\varepsilon \\[1pt]
\bm{a}_\Omega \\[1pt]
a_\vartheta \\[1pt]
\bm{b}_{\Section}
\end{pmatrix}
\in \mathbb{R}^{1+2+1+2} = \mathbb{R}^6,
\]

Let
$\bm{r}_c$ denote the centroid location. We write $\bm{r}_\perp$ for the
projection of $\bm{r}$ into the plane orthogonal to $\FrameBasis$ (so
$\bm{r}_\perp = \bm{r}$ if coordinates are chosen so that
$\FrameBasis\cdot\bm{r}=0$ on $\Omega$).

The Saint-Venant pointwise axial strain and shear strain are given (in
Ieşan's notation) by
\begin{subequations}
\begin{align}
\varepsilon(x,\bm{r};\bm{p})
&= a_\varepsilon + \bm{r}_\perp\cdot\bm{a}_\Omega
+ x\,(\bm{r}_\perp-\bm{r}_c)\cdot\bm{b}_{\Section},
\label{eq:sv-point-axial-strain} \\
\bm{\gamma}(x,\bm{r};\bm{p})
&= (a_\vartheta + c_\vartheta(\bm{b}_{\Section}))
\bigl(\FrameBasis\times \bm{r}_\perp + \nabla_\perp \omega(\bm{r}_\perp)\bigr)
\nonumber\\[-2pt]
&\quad
- \nu \bigl(\operatorname{dev}(\bm{r}_\perp\otimes\bm{r}_\perp)
- \bm{r}_\perp\otimes\bm{r}_c\bigr)\bm{b}_{\Section}
+ \bigl(\nabla_\perp \bm{\psi}(\bm{r}_\perp)\bigr)^{\!\top} \bm{b}_{\Section},
\label{eq:sv-point-shear-strain}
\end{align}
\label{eq:sv-point-strain}
\end{subequations}
where:
\begin{itemize}
\item $\omega : \Omega \to \mathbb{R}$ is the Saint-Venant torsion warping
      function,
\item $\bm{\psi} : \Omega \to \mathbb{R}^2$ is the Saint-Venant shear warping
      vector,
\item $\operatorname{dev}(\bm{r}\otimes\bm{r})$ is the deviatoric
      part of the second-order tensor $\bm{r}_\perp\otimes\bm{r}_\perp$,
\end{itemize}
The axial strain $\varepsilon$ is scalar; the shear strain $\bm{\gamma}$ is a
two-component vector tangent to the cross-section.

\paragraph{Resultants and 3D energy.}

Let $E$ and $G$ be the Young's and shear moduli, respectively. For a given
parameter vector $\bm{p}$, the stress components associated with the
Saint-Venant solution are
\[
\AxialStress = E\,\varepsilon(x,\bm{r};\bm{p}),
\qquad
\bm{\tau}_\perp = G\,\bm{\gamma}(x,\bm{r};\bm{p}).
\]

At a fixed section $x$, the classical Saint-Venant resultants associated with
$\bm{p}$ are:
\begin{itemize}
\item Axial force
\[
N(\bm{p})
=
\int_\Omega \AxialStress\,\dA
=
E\int_\Omega \varepsilon(x,\bm{r};\bm{p})\,\dA.
\]
\item Transverse shear force vector
\[
\bm{F}(\bm{p})
=
\int_\Omega \bm{\tau}_\perp\,\dA
=
G\int_\Omega \bm{\gamma}(x,\bm{r};\bm{p})\,\dA.
\]
\item Torsional moment about $\FrameBasis$
\[
T(\bm{p})
=
\int_\Omega \bigl(\bm{r}_\perp \times \bm{\tau}_\perp\bigr)\cdot\FrameBasis
\,\dA.
\]
\item Bending moment vector in the cross-sectional plane
\[
\bm{M}_\perp(\bm{p})
=
\int_\Omega \bm{r}_\perp \times 
\bigl(\AxialStress\,\FrameBasis\bigr)\,\dA.
\]
\end{itemize}
Collecting these into a resultant vector
\[
\bm{s}(\bm{p})
=
\begin{pmatrix}
N(\bm{p}) \\[2pt]
\bm{F}(\bm{p}) \\[2pt]
T(\bm{p}) \\[2pt]
\bm{M}_\perp(\bm{p})
\end{pmatrix}
\in \mathbb{R}^6,
\]
we obtain a linear map from the parameter space to resultants:
\begin{equation}
\bm{s}(\bm{p})
=
\mathbf{S}\,\bm{p},
\qquad
\mathbf{S} \in \mathbb{R}^{6\times 6}.
\label{eq:def-S-matrix}
\end{equation}
The columns of $\mathbf{S}$ are the resultants generated by the six unit
parameter states (one for each component of $\bm{p}$).

The three-dimensional elastic energy per unit length associated with a
Saint-Venant state $\bm{p}$ is
\begin{equation}
U(\bm{p})
=
\frac{1}{2}
\int_\Omega \Bigl(
E\,\varepsilon(x,\bm{r};\bm{p})^2
+
G\,\bigl|\bm{\gamma}(x,\bm{r};\bm{p})\bigr|^2
\Bigr)\,\dA.
\label{eq:sv-energy}
\end{equation}
Because $\varepsilon$ and $\bm{\gamma}$ depend linearly on $\bm{p}$, the
energy $U(\bm{p})$ is a quadratic form in $\bm{p}$:
\begin{equation}
U(\bm{p})
=
\frac{1}{2}\,\IesanVector\cdot \mathbf{K}\,\bm{p},
\qquad
\mathbf{K} \in \mathbb{R}^{6\times 6},
\label{eq:def-K-matrix}
\end{equation}
where the symmetric \emph{energy matrix} $\mathbf{K}$ has entries
\begin{equation}
K_{ij}
=
\int_\Omega \Bigl(
E\,\varepsilon(x,\bm{r};\bm{e}^{(i)})\,\varepsilon(x,\bm{r};\bm{e}^{(j)})
+
G\,\bm{\gamma}(x,\bm{r};\bm{e}^{(i)})\cdot
\bm{\gamma}(x,\bm{r};\bm{e}^{(j)})
\Bigr)\,\dA,
\label{eq:K-entries}
\end{equation}
and $\{\bm{e}^{(i)}\}_{i=1}^6$ is the canonical basis of $\mathbb{R}^6$.

\begin{remark}
The matrices $\mathbf{S}$ and $\mathbf{K}$ depend only on the cross-section
geometry, material properties, and Saint-Venant warping functions
$\omega(\bm{r}_\perp)$ and $\bm{\psi}(\bm{r}_\perp)$. They can be computed
once for a given section by solving the Saint-Venant problems for each unit
parameter state $\bm{e}^{(i)}$, evaluating the strains via
\eqref{eq:sv-point-strain}, and carrying out the cross-sectional integrals in
\eqref{eq:def-S-matrix} and \eqref{eq:K-entries}.
\end{remark}

\paragraph{General linear trace and beam strains.}

We now consider any linear trace defined at the Saint-Venant level by a
$6\times 6$ matrix $\TraceDeDpRoot$ mapping parameters to constant beam strain
coefficients:
\begin{equation}
\bm{e}_0(\bm{p})
=
\TraceDeDpRoot\,\bm{p}
\in \mathbb{R}^6,
\qquad
\bm{e}_0 =
\begin{pmatrix}
\varepsilon_0 \\[1pt]
\bm{\gamma}_{\perp,0} \\[1pt]
\kappa_\vartheta \\[1pt]
\bm{\kappa}_{\perp,0}
\end{pmatrix},
\label{eq:def-e0}
\end{equation}
where
\begin{itemize}
\item $\varepsilon_0$ is the axial strain measure,
\item $\bm{\gamma}_{\perp,0} \in \mathbb{R}^2$ is the transverse shear strain
      measure,
\item $\kappa_\vartheta$ is the twist curvature,
\item $\bm{\kappa}_{\perp,0} \in \mathbb{R}^2$ is the bending curvature
      vector.
\end{itemize}
Given $\TraceDeDpRoot$, the corresponding one-dimensional constitutive law takes
the form
\begin{equation}
\bm{s} = \mathbf{C}\,\bm{e}_0,
\qquad
\mathbf{C} = \mathbf{S}\,\TraceDpDeRoot.
\label{eq:def-C-from-T0}
\end{equation}
In general $\mathbf{C}$ need not be symmetric, and the beam strain measures
$\bm{e}_0$ need not be work-conjugate to the Saint-Venant resultants
$\bm{s}$.

\paragraph{Energetic trace: definition and construction.}

We now define an \emph{energetic trace} as one for which all six components of
$\bm{e}_0$ are work-conjugate to the classical resultants.

\begin{definition}[Energetic trace]
A linear trace with matrix $\TraceDeDpRoot$ is called \emph{energetic} if the
following two conditions hold for all Saint-Venant states $\bm{p}$ and all
variations $\delta\bm{p}$:
\begin{enumerate}
  \item The one-dimensional internal virtual work matches the three-dimensional
        virtual work on the Saint-Venant class:
        \begin{equation}
        \bm{s}(\bm{p})\cdot \delta\bm{e}_0(\bm{p})
        =
        \int_\Omega \Bigl(
        \AxialStress(\bm{p})\,\delta\varepsilon(\delta\bm{p})
        +
        \bm{\tau}_\perp(\bm{p})\cdot\delta\bm{\gamma}(\delta\bm{p})
        \Bigr)\,\dA.
        \label{eq:ener-virtual-work}
        \end{equation}
  \item The one-dimensional strain energy density
        \begin{equation}
        U_{\mathrm{1D}}(\bm{p})
        :=
        \frac{1}{2}\,\bm{s}(\bm{p})\cdot \bm{e}_0(\bm{p})
        \end{equation}
        coincides with the three-dimensional Saint-Venant energy
        $U(\bm{p})$ in \eqref{eq:sv-energy}:
        \begin{equation}
        U_{\mathrm{1D}}(\bm{p}) = U(\bm{p})
        \quad\text{for all }\bm{p}\in\mathbb{R}^6.
        \label{eq:ener-energy-equality}
        \end{equation}
\end{enumerate}
\end{definition}

We now show that these requirements uniquely determine the energetic trace
matrix $\TraceDeDpRoot$ in terms of the matrices $\mathbf{S}$ and $\mathbf{K}$.

\begin{proposition}[Energetic trace matrix]
\label{prop:energetic-T0}
Assume that the resultant matrix $\mathbf{S}\in\mathbb{R}^{6\times 6}$ and the
energy matrix $\mathbf{K}\in\mathbb{R}^{6\times 6}$ defined above are both
invertible. 
Then there exists a unique energetic trace, and its trace matrix
$\TraceDeDpRoot$ is given by
\begin{equation}
\TraceDeDpRoot
=
\mathbf{S}^{-\top}\,\mathbf{K}.
\label{eq:T0-energetic}
\end{equation}
\end{proposition}

\begin{proof}
Let $\bm{p}$ and $\delta\bm{p}$ be arbitrary parameter vectors. On the one
hand, the three-dimensional virtual work associated with a Saint-Venant state
$\bm{p}$ and variation $\delta\bm{p}$ is
\begin{equation*}
\delta W_{\mathrm{3D}}
=
\int_\Omega \Bigl(
\AxialStress(\bm{p})\,\delta\varepsilon(\delta\bm{p})
+
\bm{\tau}_\perp(\bm{p})\cdot\delta\bm{\gamma}(\delta\bm{p})
\Bigr)\,\dA.
\end{equation*}
Because the strains depend linearly on $\bm{p}$, this virtual work is a
bilinear form in $(\bm{p},\delta\bm{p})$. By differentiating
\eqref{eq:sv-energy} we obtain
\begin{equation}
\delta W_{\mathrm{3D}}
=
\frac{\partial}{\partial t}
U(\bm{p}+t\,\delta\bm{p})\big|_{t=0}
=
\IesanVector\cdot \mathbf{K}\,\delta\bm{p}.
\label{eq:3D-work-K}
\end{equation}

On the other hand, the one-dimensional internal virtual work is
\begin{equation}
\delta W_{\mathrm{1D}}
=
\bm{s}(\bm{p})\cdot \delta\bm{e}_0(\bm{p})
=
\bigl(\mathbf{S}\bm{p}\bigr)^\top
\TraceDeDpRoot\,\delta\bm{p}
=
\IesanVector\cdot \mathbf{S}^\top \TraceDeDpRoot\,\delta\bm{p}.
\label{eq:1D-work-T0}
\end{equation}
Requiring $\delta W_{\mathrm{1D}}=\delta W_{\mathrm{3D}}$ for all
$\bm{p},\delta\bm{p}$ yields
\begin{equation}
\mathbf{S}^\top \TraceDeDpRoot = \mathbf{K}.
\label{eq:work-equality-matrix}
\end{equation}
Since $\mathbf{S}$ is invertible by assumption, this gives
\begin{equation}
\TraceDeDpRoot = \mathbf{S}^{-\top}\,\mathbf{K},
\end{equation}
which is unique.

We now check the energy equality \eqref{eq:ener-energy-equality}. With
$\bm{e}_0 = \TraceDeDpRoot\bm{p}$ and $\bm{s}=\mathbf{S}\bm{p}$, we have
\begin{equation*}
\bm{s}\cdot\bm{e}_0
=
(\mathbf{S}\bm{p})^\top(\TraceDeDpRoot\bm{p})
=
\IesanVector\cdot \mathbf{S}^\top \TraceDeDpRoot\,\bm{p}
=
\IesanVector\cdot \mathbf{K}\,\bm{p},
\end{equation*}
using \eqref{eq:work-equality-matrix}. Hence
\begin{equation*}
U_{\mathrm{1D}}(\bm{p})
=
\frac{1}{2}\,\bm{s}\cdot\bm{e}_0
=
\frac{1}{2}\,\IesanVector\cdot \mathbf{K}\,\bm{p}
=
U(\bm{p}),
\end{equation*}
which establishes \eqref{eq:ener-energy-equality}. Thus the trace with
$\TraceDeDpRoot = \mathbf{S}^{-\top}\mathbf{K}$ is energetic in the sense of the
definition, and it is unique.
\end{proof}

\begin{remark}[Full work conjugacy]
The construction \eqref{eq:T0-energetic} makes \emph{all} six components of
$\bm{e}_0$ work-conjugate to the corresponding components of the resultant
vector $\bm{s}$. Indeed, by \eqref{eq:3D-work-K} and
\eqref{eq:1D-work-T0}, we have for any state $\bm{p}$ and variation
$\delta\bm{p}$
\[
\bm{s}(\bm{p})\cdot\delta\bm{e}_0(\bm{p})
=
\IesanVector\cdot \mathbf{K}\,\delta\bm{p}
=
\int_\Omega \Bigl(
\AxialStress(\bm{p})\,\delta\varepsilon(\delta\bm{p})
+
\bm{\tau}_\perp(\bm{p})\cdot\delta\bm{\gamma}(\delta\bm{p})
\Bigr)\,\dA.
\]
Thus the axial strain $\varepsilon_0$, shear strain $\bm{\gamma}_{\perp,0}$,
twist curvature $\kappa_\vartheta$, and bending curvature $\bm{\kappa}_{\perp,0}$
are each energy-conjugate to the corresponding resultants
$(N,\bm{F},T,\bm{M}_\perp)$ in the classical Saint-Venant sense.
\end{remark}

\paragraph{Energetic constitutive matrix.}

Given $\TraceDeDpRoot = \mathbf{S}^{-\top}\mathbf{K}$, the one-dimensional
constitutive matrix associated with the energetic trace is
\begin{equation*}
\mathbf{C}
=
\mathbf{S}\,\TraceDpDeRoot
=
\mathbf{S}\,\mathbf{K}^{-1}\mathbf{S}^\top.
% \label{eq:C-energetic}
\end{equation*}
This matrix is symmetric and positive definite (since $\mathbf{K}$ is symmetric
positive definite and $\mathbf{S}$ is invertible), and the energy can be
written equivalently as
\begin{equation}
U(\bm{p})
=
\frac{1}{2}\,\IesanVector\cdot\mathbf{K}\,\bm{p}
=
\frac{1}{2}\,\bm{e}_0^\top \mathbf{C}\,\bm{e}_0
=
\frac{1}{2}\,\bm{s}^\top \mathbf{C}^{-1}\bm{s}.
\end{equation}

\begin{remark}[Shear and torsion blocks]
Although the energetic construction \eqref{eq:T0-energetic} acts on the full
six-dimensional parameter space, it is often useful to inspect the shear and
torsion blocks explicitly. By partitioning the parameter vector as
$\bm{p}=(a_\varepsilon,\bm{a}_\Omega,a_\vartheta,\bm{b}_{\Section})^\top$ and the
resultants as $\bm{s}=(N,\bm{F},T,\bm{M}_\perp)^\top$, the matrices
$\mathbf{S}$ and $\mathbf{K}$ can be written in block form, and the energetic
trace matrix $\TraceDeDpRoot = \mathbf{S}^{-\top}\mathbf{K}$ induces:
\begin{itemize}
\item an axial block $(N,\varepsilon_0)$,
\item a $2\times 2$ shear block relating $\bm{F}$ and $\bm{\gamma}_{\perp,0}$,
\item a torsional block relating $T$ and $\kappa_\vartheta$,
\item a bending block relating $\bm{M}_\perp$ and $\bm{\kappa}_{\perp,0}$.
\end{itemize}
In particular, the shear and torsion blocks of $\mathbf{C}$ coincide with the
classical energy-conjugate shear stiffness and torsional stiffness obtained by
integrating the Saint-Venant stress fields generated by
\eqref{eq:sv-point-strain}.
\end{remark}

\begin{remark}[Relation to other traces and the trace matrices $T_1$]
Within the Saint-Venant solution class, the beam strain fields
$\bm{e}(x;\bm{p})$ are at most affine in $x$:
\[
\bm{e}(x;\bm{p}) = \bm{e}_0(\bm{p}) + x\,\bm{e}_1(\bm{p}),
\]
where the gradient coefficients $\bm{e}_1(\bm{p})$ are independent of the
trace choice. 
Different traces (pointwise/Timoshenko, section-average, energetic)
correspond to different choices of the constant matrix $\TraceDeDpRoot$ mapping
$\bm{p}$ to $\bm{e}_0(\bm{p})$, while the first-order trace matrix
$\TraceDeDpOne$ (mapping $\bm{p}$ to $\bm{e}_1(\bm{p})$) is universal in the
Saint-Venant class. 
The energetic trace is the unique choice of $\TraceDeDpRoot$
for which the resulting one-dimensional strains are fully work-conjugate to the
classical resultants.
\end{remark}


% \subsubsection{The Energetic (Renton-Type) Trace}

% We now introduce the \emph{energetic} trace, which is characterized by the 
% requirement that the one-dimensional strain measures be \emph{work-conjugate} 
% to the Saint-Venant resultants. The key adjustment, relative to the pointwise 
% or section-average traces, is in the definition of the transverse shear strain: 
% it is chosen so that its work with the shear resultant reproduces the exact 
% Saint-Venant shear energy.

% \paragraph{Problem I / Problem II decomposition.}

% Recall that the Saint-Venant displacement is decomposed into the sum of two 
% solution families:
% \begin{equation}
% \hat{\bm{u}}
% =
% \hat{\bm{u}}^{\rm I}(\bm{a})
% +
% \hat{\bm{u}}^{\rm I\!I}(\bm{b}),
% \qquad
% \bm{a} = (a_\varepsilon,\bm{a}_\Omega,a_\vartheta) \in \mathbb{R}^4,
% \quad
% \bm{b} = \bm{b}_{\Section} \in \mathbb{R}^2.
% \end{equation}

% For any Saint-Venant state, the resultant vector
% \begin{equation}
% \bm{s}
% =
% \begin{pmatrix}
% N \\ T \\ \bm{F} \\ \FrameBasis \times \bm{M}
% \end{pmatrix}
% \in \mathbb{R}^6
% \end{equation}
% depends linearly on $\bm{p}$ and can be written as
% \begin{equation}
% \bm{s}
% =
% \bm{s}^{\rm I}(\bm{a})
% +
% \bm{s}^{\rm I\!I}(\bm{b}),
% \qquad
% \bm{s}^{\rm I}(\bm{a}) = \mathbf{R}^{\rm I}\,\bm{a},
% \quad
% \bm{s}^{\rm I\!I}(\bm{b}) = \mathbf{R}^{\rm I\!I}\,\bm{b},
% \end{equation}
% where $\mathbf{R}^{\rm I} \in \mathbb{R}^{6\times 4}$ and 
% $\mathbf{R}^{\rm I\!I} \in \mathbb{R}^{6\times 2}$ are the 
% Problem~$\mathrm{I}$ and Problem~$\mathrm{II}$ resultant matrices. The 
% transverse shear resultant is the subvector
% \begin{equation}
% \bm{F}
% % =
% % \bm{F}^{\rm I}(\bm{a}) + \bm{F}^{\rm I\!I}(\bm{b})
% =
% \BishearTensorRB\,\bm{b},
% \end{equation}
% where 
% $\BishearTensor^{\rm I} \in \mathbb{R}^{2\times 4}$ and 
% $\BishearTensor^{\rm I\!I} \in \mathbb{R}^{2\times 2}$ collect the shear 
% resultants produced by the two problems (``bishear'' tensors).

% The Saint-Venant shear energy per unit length is quadratic in $(\bm{a},\bm{b})$,
% \begin{equation}
% \mathcal{U}_\mathrm{sh}(\bm{a},\bm{b})
% =
% \frac{1}{2}
% \begin{pmatrix}
% \bm{a} \\ \bm{b}
% \end{pmatrix}^{\!\top}
% \begin{pmatrix}
% \mathbf{H}^{\rm I\!I}_{\mathrm{aa}} & \mathbf{H}^{\rm I\!I}_{\mathrm{ab}} \\
% \mathbf{H}^{\rm I\!I}_{\mathrm{ba}} & \mathbf{H}^{\rm I\!I}_{\mathrm{bb}}
% \end{pmatrix}
% \begin{pmatrix}
% \bm{a} \\ \bm{b}
% \end{pmatrix},
% \label{eq:sv-shear-energy-block}
% \end{equation}
% with a symmetric positive-definite ``shear energy tensor''
% $\mathbf{H}^{\rm I\!I} \in \mathbb{R}^{6\times 6}$ and blocks
% $\mathbf{H}^{\rm I\!I}_{\mathrm{aa}}\in\mathbb{R}^{4\times 4}$,
% $\mathbf{H}^{\rm I\!I}_{\mathrm{bb}}\in\mathbb{R}^{2\times 2}$, etc.
% In particular, the Problem~$\mathrm{I}$ shear energy is 
% $\tfrac{1}{2}\bm{a}^\top\mathbf{H}^{\rm I\!I}_{\mathrm{aa}}\bm{a}$, and the 
% Problem~$\mathrm{II}$ shear energy is 
% $\tfrac{1}{2}\bm{b}^\top\mathbf{H}^{\rm I\!I}_{\mathrm{bb}}\bm{b}$.


% \begin{definition}[Energetic trace]
% \label{def:energetic-trace}
% An \emph{energetic trace} is a Saint-Venant linear trace $\mathcal{T}_{\mathrm{en}}$ 
% such that:
% \begin{enumerate}
%   \item It preserves the canonical identification for axial strain, torsion, 
%   and bending curvature:
%   \begin{equation}
%   \varepsilon_0 = a_\varepsilon, 
%   \qquad
%   \bm{\kappa}_{\perp,0} = -\FrameBasis \times \bm{a}_\Omega,
%   \qquad
%   \varkappa_0 = a_\vartheta + \partial_{\bm{b}} c_\vartheta \cdot \bm{b}_{\Section}.
%   \label{eq:ener-canonical-non-shear}
%   \end{equation}
%   \item Its transverse shear strain $\bm{\gamma}_{\perp,0}$ is 
%   \emph{work-conjugate} to the Saint-Venant shear resultant in the sense that 
%   the exact Saint-Venant shear energy is represented as
%   \begin{equation}
%   \mathcal{U}_\mathrm{sh}(\bm{a},\bm{b})
%   =
%   \frac{1}{2}\,\bm{\gamma}_{\perp,0}(\bm{a},\bm{b})^\top 
%   \mathbf{H}_{\mathrm{sh}}
%   \,\bm{\gamma}_{\perp,0}(\bm{a},\bm{b}),
%   \label{eq:ener-shear-energy}
%   \end{equation}
%   for some symmetric positive-definite \emph{shear stiffness} 
%   $\mathbf{H}_{\mathrm{sh}}\in\mathbb{R}^{2\times 2}$, and
%   \begin{equation}
%   \bm{F}(\bm{a},\bm{b})
%   =
%   \mathbf{H}_{\mathrm{sh}}\,\bm{\gamma}_{\perp,0}(\bm{a},\bm{b})
%   \label{eq:ener-shear-work-conj}
%   \end{equation}
%   for all Saint-Venant states $(\bm{a},\bm{b})$.
% \end{enumerate}
% \end{definition}

% The first condition simply fixes the gauge for non-shear modes to match the 
% standard Cosserat interpretation. The second condition enforces exact equality 
% of three-dimensional and one-dimensional shear work and energy, making 
% $\bm{\gamma}_{\perp,0}$ the \emph{energetic shear measure}.

% \paragraph{Algebraic structure of the shear block.}

% By linearity on the Saint-Venant parameter space, the shear strain must have 
% the form
% \begin{equation}
% \bm{\gamma}_{\perp,0}(\bm{a},\bm{b})
% =
% \mathbf{X}^{\rm I}\,\bm{a}
% +
% \mathbf{X}^{\rm I\!I}\,\bm{b},
% \qquad
% \mathbf{X}^{\rm I} \in \mathbb{R}^{2\times 4},
% \quad
% \mathbf{X}^{\rm I\!I} \in \mathbb{R}^{2\times 2}.
% \end{equation}
% Similarly, the transverse shear resultant can be written as
% \begin{equation}
% \bm{F}(\bm{a},\bm{b})
% =
% \BishearTensor^{\rm I}\,\bm{a}
% +
% \BishearTensor^{\rm I\!I}\,\bm{b},
% \end{equation}
% with the bishear tensors introduced above.

% Substituting these expansions into the work-conjugacy requirement 
% \eqref{eq:ener-shear-work-conj} gives, for all $(\bm{a},\bm{b})$,
% \begin{equation}
% \BishearTensor^{\rm I}\,\bm{a}
% +
% \BishearTensor^{\rm I\!I}\,\bm{b}
% =
% \mathbf{H}_{\mathrm{sh}}
% \big(\mathbf{X}^{\rm I}\,\bm{a} + \mathbf{X}^{\rm I\!I}\,\bm{b}\big),
% \end{equation}
% or equivalently,
% \begin{equation}
% \BishearTensor^{\rm I} = \mathbf{H}_{\mathrm{sh}}\,\mathbf{X}^{\rm I},
% \qquad
% \BishearTensor^{\rm I\!I} = \mathbf{H}_{\mathrm{sh}}\,\mathbf{X}^{\rm I\!I}.
% \label{eq:X-from-BH}
% \end{equation}

% On the other hand, equating the shear energy expressions 
% \eqref{eq:sv-shear-energy-block} and \eqref{eq:ener-shear-energy} gives
% \begin{equation}
% \frac{1}{2}
% \begin{pmatrix}
% \bm{a} \\ \bm{b}
% \end{pmatrix}^{\!\top}
% \begin{pmatrix}
% \mathbf{H}^{\rm I\!I}_{\mathrm{aa}} & \mathbf{H}^{\rm I\!I}_{\mathrm{ab}} \\
% \mathbf{H}^{\rm I\!I}_{\mathrm{ba}} & \mathbf{H}^{\rm I\!I}_{\mathrm{bb}}
% \end{pmatrix}
% \begin{pmatrix}
% \bm{a} \\ \bm{b}
% \end{pmatrix}
% =
% \frac{1}{2}
% \big(\mathbf{X}^{\rm I}\bm{a} + \mathbf{X}^{\rm I\!I}\bm{b}\big)^\top
% \mathbf{H}_{\mathrm{sh}}
% \big(\mathbf{X}^{\rm I}\bm{a} + \mathbf{X}^{\rm I\!I}\bm{b}\big),
% \end{equation}
% for all $(\bm{a},\bm{b})$. Thus
% \begin{equation}
% \begin{pmatrix}
% \mathbf{H}^{\rm I\!I}_{\mathrm{aa}} & \mathbf{H}^{\rm I\!I}_{\mathrm{ab}} \\
% \mathbf{H}^{\rm I\!I}_{\mathrm{ba}} & \mathbf{H}^{\rm I\!I}_{\mathrm{bb}}
% \end{pmatrix}
% =
% \begin{pmatrix}
% \mathbf{X}^{\mathrm{I}\top} \\ \mathbf{X}^{\mathrm{II}\top}
% \end{pmatrix}
% \mathbf{H}_{\mathrm{sh}}
% \begin{pmatrix}
% \mathbf{X}^{\rm I} & \mathbf{X}^{\rm I\!I}
% \end{pmatrix}.
% \label{eq:H-block-factorization}
% \end{equation}

% Equations \eqref{eq:X-from-BH}--\eqref{eq:H-block-factorization} determine 
% $\mathbf{X}^{\rm I}$ and $\mathbf{X}^{\rm I\!I}$ (and hence the shear 
% rows of the trace matrix) in terms of the Saint-Venant bishear tensors and the 
% shear energy tensor. In particular, if $\mathbf{H}_{\mathrm{sh}}$ is chosen as
% \begin{equation}
% \mathbf{H}_{\mathrm{sh}} = \mathbf{H}^{\rm I\!I}_{\mathrm{bb}},
% \end{equation}
% then the Problem~$\mathrm{II}$ block simplifies to
% \begin{equation}
% \mathbf{X}^{\rm I\!I} = \mathbf{H}_{\mathrm{sh}}^{-1}\,\BishearTensor^{\rm I\!I},
% \end{equation}
% and the Problem~$\mathrm{I}$ block follows from 
% \eqref{eq:H-block-factorization} as the unique solution of
% \begin{equation}
% \mathbf{H}^{\rm I\!I}_{\mathrm{aa}} 
% =
% \mathbf{X}^{\mathrm{I}\top}\mathbf{H}_{\mathrm{sh}}\mathbf{X}^{\rm I}, 
% \qquad
% \mathbf{H}^{\rm I\!I}_{\mathrm{ab}}
% =
% \mathbf{X}^{\mathrm{I}\top}\mathbf{H}_{\mathrm{sh}}\mathbf{X}^{\rm I\!I}.
% \end{equation}
% In practice, and following the usual Ieşan notation, one may identify
% \begin{equation}
% \mathbf{H}_{\mathrm{sh}} \equiv \mathbf{H},
% \qquad
% \BishearTensor^{\rm I} \equiv \BishearTensorRB,
% \end{equation}
% leading to the compact expression
% \begin{equation}
% \mathbf{X}^{\rm I} = \mathbf{H}^{-\top}\,\BishearTensorRB^\top,
% \qquad
% \mathbf{X}^{\rm I\!I} = \mathbf{0},
% \label{eq:X-ener-concrete}
% \end{equation}
% for the energetic shear measure used in the classical Renton-type construction.

% \paragraph{Energetic trace matrices.}

% With the above notation, the energetic trace acts on the Ieşan parameters as
% \begin{equation}
% \bm{e}_{0}(\bm{p})
% =
% \TraceDeDpRoot\,\bm{p},
% \qquad
% \bm{p} = (a_\varepsilon,\bm{a}_\Omega,a_\vartheta,\bm{b}_{\Section})^\top,
% \end{equation}
% with
% \begin{equation}
% \TraceDeDpRoot =
% \begin{pmatrix}
% \FrameBasis^\top & \mathbf{0}^\top & 0 & \mathbf{0}^\top \\[4pt]
% \mathbf{0} & \mathbf{X}^{\rm I} & \mathbf{0} & \mathbf{X}^{\rm I\!I} \\[4pt]
% 0 & \mathbf{0}^\top & 1 & \partial_{\bm{b}} c_\vartheta \\[4pt]
% \mathbf{0} & -\FrameBasis^{\times} & \mathbf{0} & \mathbf{0}
% \end{pmatrix},
% \label{eq:T0-ener-general}
% \end{equation}
% where $\mathbf{X}^{\rm I},\mathbf{X}^{\rm I\!I}$ are determined by 
% \eqref{eq:X-from-BH}--\eqref{eq:H-block-factorization}. In the classical 
% Renton-type energetic shear gauge, \eqref{eq:X-ener-concrete} gives the 
% concrete form
% \begin{equation}
% \TraceDeDpRoot =
% \begin{pmatrix}
% \FrameBasis^\top & \mathbf{0}^\top & 0 & \mathbf{0}^\top \\[4pt]
% \mathbf{0} & \mathbf{H}^{-\top}\,\BishearTensorRB^\top & \mathbf{0} & \mathbf{0} \\[4pt]
% 0 & \mathbf{0}^\top & 1 & \partial_{\bm{b}} c_\vartheta \\[4pt]
% \mathbf{0} & -\FrameBasis^{\times} & \mathbf{0} & \mathbf{0}
% \end{pmatrix}.
% \label{eq:T0-ener-Renton}
% \end{equation}

% \begin{remark}[Affine-in-$x$ structure and $T_{1}$]
% Within the Saint-Venant class, all traces considered in this work (pointwise, 
% section-average, energetic/Renton) produce beam strains that are at most linear 
% in $x$. The \emph{gradient} coefficients $(\bm{\gamma}_1,\bm{\kappa}_1)$ are 
% universal: they depend only on the flexural parameters $\bm{b}_{\Section}$ and are 
% independent of the trace gauge. Consequently, the first-order trace matrix 
% $\TraceDeDpOne$ coincides with that of any other trace in this 
% SV-equivalent family (e.g.\ the section-average or pointwise trace), and the 
% distinction between traces resides entirely in the constant matrix 
% $\TraceDeDpRoot$.
% \end{remark}

% \begin{remark}[Problem~$\mathrm{I}$ vs.\ Problem~$\mathrm{II}$ blocks]
% It is often convenient to write
% \begin{equation}
% \TraceDeDpRoot =
% \big(\,\mathbf{T}^{\rm I}_{0} \;\big|\; \mathbf{T}^{\rm I\!I}_{0}\,\big),
% \end{equation}
% with $\mathbf{T}^{\rm I}_{0} \in \mathbb{R}^{6\times 4}$ 
% acting on the Problem~$\mathrm{I}$ parameters $(a_\varepsilon,\bm{a}_\Omega,a_\vartheta)$ 
% and $\mathbf{T}^{\rm I\!I}_{0} \in \mathbb{R}^{6\times 2}$ acting 
% on the Problem~$\mathrm{II}$ parameters $\bm{b}_{\Section}$. In the Renton-type 
% gauge \eqref{eq:T0-ener-Renton}, Problem~$\mathrm{II}$ contributes only to the 
% twist curvature via $\partial_{\bm{b}}c_\vartheta$, while the shear gauge is 
% fixed entirely by Problem~$\mathrm{I}$ through the energetic condition 
% \eqref{eq:X-ener-concrete}.
% \end{remark}

\subsubsection{The Energy-Conjugate (Renton-Romano) Trace}

The energy-conjugate trace is defined by requiring that the beam strains 
$(\bm{\gamma}, \bm{\kappa})$ be work-conjugate to the classical resultants 
$(\bm{n}, \bm{m})$. That is, for any variation $\delta\hat{\bm{u}}$ within the 
Saint-Venant class:
\begin{equation}
\delta\bm{\gamma} \cdot \bm{n} + \delta\bm{\kappa} \cdot \bm{m} 
= \int_\Section \delta\bm{\varepsilon} : \bm{\sigma} \, \dA.
\label{eq:energy-conjugacy}
\end{equation}

We derive this trace by examining the structure of the Saint-Venant strain 
field and identifying which strain measures satisfy \eqref{eq:energy-conjugacy}.

\paragraph{Saint-Venant strain fields.}

From \eqref{eq:iesan-point-strain}, the pointwise strains in a Saint-Venant 
solution are:
\begin{align}
\PointAxialStrain &= a_\varepsilon + \bm{r} \cdot \bm{a}_\Omega 
+ x\,(\bm{r} - \CentroidLocation) \cdot \bm{b}_{\Section}, 
\label{eq:sv-axial-strain} \\
\PointShearStrain &= (a_\vartheta + c_\vartheta)(\FrameBasis \times \bm{r} + \nabla\WarpTwistIesan) 
+ \bm{W}_{(b)}\,\bm{b}_{\Section} + (\nabla\bm{\WarpShearIesan})^\top\,\bm{b}_{\Section},
\label{eq:sv-shear-strain}
\end{align}

The axial force, shear force, torque, and bending moment are:
\begin{align*}
N &= \int_\Section \sigma \, \dA, \\
\bm{F} &= \int_\Section \bm{\tau} \, \dA, \\
T &= \int_\Section \bm{r} \times \bm{\tau} \cdot \FrameBasis \, \dA 
= \int_\Section (\FrameBasis \times \bm{r}) \cdot \bm{\tau} \, \dA, \\
\bm{M} &= \int_\Section \bm{r} \times (\sigma\,\FrameBasis) \, \dA 
= \FrameBasis \times \int_\Section \bm{r}\,\sigma \, \dA.
\end{align*}

For a linear elastic material with Young's modulus $E$ and shear modulus $G$:
\begin{align}
\sigma &= E\,\PointAxialStrain, \\
\bm{\tau} &= G\,\PointShearStrain.
\end{align}

\paragraph{Work-conjugacy conditions.}

The right-hand side of \eqref{eq:energy-conjugacy} is:
\begin{equation}
\int_\Section \delta\bm{\varepsilon} : \bm{\sigma} \, \dA 
= \int_\Section \delta\PointAxialStrain\,\sigma \, \dA 
+ \int_\Section \delta\PointShearStrain \cdot \bm{\tau} \, \dA.
\label{eq:rhs-energy}
\end{equation}

We seek beam strains $(\epsilon, \bm{\gamma}_\perp, \varkappa, \bm{\kappa}_\perp)$ 
such that:
\begin{equation}
\delta\epsilon\,N + \delta\bm{\gamma}_\perp \cdot \bm{F} 
+ \delta\varkappa\,T + \delta\bm{\kappa}_\perp \cdot \bm{M} 
= \int_\Section \delta\PointAxialStrain\,\sigma \, \dA 
+ \int_\Section \delta\PointShearStrain \cdot \bm{\tau} \, \dA.
\label{eq:conjugacy-expanded}
\end{equation}

\subparagraph{Axial strain $\epsilon$.}

Consider a variation with $\delta a_\varepsilon \neq 0$ and all other parameters zero. Then:
\begin{equation*}
\delta\PointAxialStrain = \delta a_\varepsilon, \qquad \delta\PointShearStrain = \mathbf{0}.
\end{equation*}
The right-hand side of \eqref{eq:conjugacy-expanded} becomes:
\begin{equation*}
\int_\Section \delta a_\varepsilon\,\sigma \, \dA = \delta a_\varepsilon\,N.
\end{equation*}
For this to equal $\delta\epsilon\,N$, we require:
\begin{equation*}
\epsilon^{(\varepsilon)} = a_\varepsilon.
% \label{eq:conj-axial-ext}
\end{equation*}

\subparagraph{Curvature $\bm{\kappa}_\perp$.}

Consider a variation with $\delta\bm{a}_\Omega \neq \mathbf{0}$ and all other parameters zero. Then:
\begin{equation*}
\delta\PointAxialStrain = \bm{r} \cdot \delta\bm{a}_\Omega, \qquad \delta\PointShearStrain = \mathbf{0}.
\end{equation*}
The right-hand side becomes:
\begin{equation*}
\int_\Section (\bm{r} \cdot \delta\bm{a}_\Omega)\,\sigma \, \dA 
= \delta\bm{a}_\Omega \cdot \int_\Section \bm{r}\,\sigma \, \dA.
\end{equation*}
Now, $\bm{M} = \FrameBasis \times \int_\Section \bm{r}\,\sigma \, \dA$, so 
$\int_\Section \bm{r}\,\sigma \, \dA = -\FrameBasis \times \bm{M}$.

For the left-hand side to match, we need:
\begin{equation*}
\delta\bm{\kappa}_\perp \cdot \bm{M} = -\delta\bm{a}_\Omega \cdot (\FrameBasis \times \bm{M}) 
= (\FrameBasis \times \delta\bm{a}_\Omega) \cdot \bm{M}.
\end{equation*}
Thus:
\begin{equation*}
\bm{\kappa}_\perp^{(a)} = \FrameBasis \times \bm{a}_\Omega = -\FrameBasis^\times\,\bm{a}_\Omega.
\end{equation*}

\subparagraph{Twist $\varkappa$.}

Consider a variation with $\delta a_\vartheta \neq 0$ and all other parameters zero. Then:
\begin{equation*}
\delta\PointAxialStrain = 0, \qquad 
\delta\PointShearStrain = \delta a_\vartheta\,(\FrameBasis \times \bm{r} + \nabla\WarpTwistIesan).
\end{equation*}
The right-hand side becomes:
\begin{equation*}
\int_\Section \delta a_\vartheta\,(\FrameBasis \times \bm{r} + \nabla\WarpTwistIesan) \cdot \bm{\tau} \, \dA 
= \delta a_\vartheta \int_\Section (\FrameBasis \times \bm{r}) \cdot \bm{\tau} \, \dA 
+ \delta a_\vartheta \int_\Section \nabla\WarpTwistIesan \cdot \bm{\tau} \, \dA.
\end{equation*}

The first integral is $T$ by definition. For the second integral, using the 
divergence theorem and the traction-free boundary condition on the lateral surface:
\begin{equation*}
\int_\Section \nabla\WarpTwistIesan \cdot \bm{\tau} \, \dA 
= \oint_{\partial\Section} \WarpTwistIesan\,(\bm{\tau} \cdot \bm{\nu}) \, \mathrm{d}s 
- \int_\Section \WarpTwistIesan\,(\nabla \cdot \bm{\tau}) \, \dA = 0,
\end{equation*}
since $\bm{\tau} \cdot \bm{\nu} = 0$ on $\partial\Section$ and $\nabla \cdot \bm{\tau} = 0$ 
in equilibrium.

Thus the right-hand side is simply $\delta a_\vartheta\,T$, and we require:
\begin{equation}
\varkappa^{(\vartheta)} = a_\vartheta.
\label{eq:conj-twist-torsion}
\end{equation}

\subparagraph{Shear strain $\bm{\gamma}_\perp$.}

Consider a variation with $\delta\bm{b}_{\Section} \neq \mathbf{0}$ and all other parameters zero. 
Then:
\begin{align*}
\delta\PointAxialStrain &= x\,(\bm{r} - \CentroidLocation) \cdot \delta\bm{b}_{\Section}, \\
\delta\PointShearStrain &= \delta c_\vartheta\,(\FrameBasis \times \bm{r} + \nabla\WarpTwistIesan) 
+ \bm{W}_{(b)}\,\delta\bm{b}_{\Section} + (\nabla\bm{\WarpShearIesan})^\top\,\delta\bm{b}_{\Section},
\end{align*}
where $\delta c_\vartheta = \CenterShearLocation \cdot \delta\bm{b}_{\Section}$.

The axial contribution to the right-hand side is:
\begin{equation*}
\int_\Section x\,(\bm{r} - \CentroidLocation) \cdot \delta\bm{b}_{\Section}\,\sigma \, \dA 
= x\,\delta\bm{b}_{\Section} \cdot \int_\Section (\bm{r} - \CentroidLocation)\,\sigma \, \dA.
\end{equation*}
This is $x$-dependent and involves $\FrameBasis \times \bm{M}$ (shifted to centroid). 
This will contribute to $\delta\bm{\kappa}_\perp \cdot \bm{M}$ terms.

The shear contribution is:
\begin{align*}
&\int_\Section \delta\PointShearStrain \cdot \bm{\tau} \, \dA \\
&= \delta c_\vartheta \int_\Section (\FrameBasis \times \bm{r} + \nabla\WarpTwistIesan) \cdot \bm{\tau} \, \dA 
+ \delta\bm{b}_{\Section} \cdot \int_\Section \bm{W}_{(b)}^\top\,\bm{\tau} \, \dA 
+ \delta\bm{b}_{\Section} \cdot \int_\Section \nabla\bm{\WarpShearIesan}\,\bm{\tau} \, \dA.
\end{align*}

The first integral gives $\delta c_\vartheta\,T = (\CenterShearLocation \cdot \delta\bm{b}_{\Section})\,T$ 
(contributing to twist-shear coupling).

For the shear terms, we use integration by parts. Consider:
\begin{equation*}
\int_\Section \nabla\bm{\WarpShearIesan}\,\bm{\tau} \, \dA 
= \oint_{\partial\Section} \bm{\WarpShearIesan}\,(\bm{\tau} \cdot \bm{\nu}) \, \mathrm{d}s 
- \int_\Section \bm{\WarpShearIesan}\,(\nabla \cdot \bm{\tau}) \, \dA = \mathbf{0},
\end{equation*}
by the same argument as before.

For the Poisson term:
\begin{equation*}
\int_\Section \bm{W}_{(b)}^\top\,\bm{\tau} \, \dA 
= -\nu \int_\Section \big(\mathrm{dev}(\bm{r} \otimes \bm{r}) 
- \bm{r} \otimes \CentroidLocation\big)^\top\,\bm{\tau} \, \dA.
\end{equation*}

Define the \emph{bishear tensor}:
\begin{equation}
\mathbf{H}_W := \int_\Section \bm{W}_{(b)}^\top\,\bm{\tau} \, \dA 
= \int_\Section \bm{W}_{(b)}^\top\,G\,\PointShearStrain \, \dA.
\label{eq:def-bishear-W}
\end{equation}

This is a $2 \times 2$ tensor (mapping $\bm{b}_{\Section}$ to a transverse vector).

The shear force is $\bm{F} = \int_\Section \bm{\tau} \, \dA$. We want 
$\delta\bm{\gamma}_\perp \cdot \bm{F}$ to capture the shear work. 

From the Saint-Venant shear stress distribution under pure flexure 
($\bm{b}_{\Section}$ loading), the shear force is:
\begin{equation}
\bm{F} = G \int_\Section \PointShearStrain \, \dA 
= G \int_\Section \big( c_\vartheta\,(\FrameBasis \times \bm{r} + \nabla\WarpTwistIesan) 
+ \bm{W}_{(b)} + (\nabla\bm{\WarpShearIesan})^\top \big)\,\bm{b}_{\Section} \, \dA.
\end{equation}

Using $\int_\Section (\FrameBasis \times \bm{r}) \, \dA = \FrameBasis \times (A\,\CentroidLocation)$, 
$\int_\Section \nabla\WarpTwistIesan \, \dA = \mathbf{0}$ (by divergence theorem), 
$\int_\Section \nabla\bm{\WarpShearIesan} \, \dA = \mathbf{0}$:
\begin{equation}
\bm{F} = G \big( c_\vartheta\,A\,(\FrameBasis \times \CentroidLocation) 
+ \SectionAverage{\bm{W}_{(b)}}\,A \big)\,\bm{b}_{\Section}.
\end{equation}

For energy conjugacy, we need the shear strain $\bm{\gamma}_\perp$ such that:
\begin{equation}
\delta\bm{\gamma}_\perp \cdot \bm{F} = \int_\Section \delta\PointShearStrain \cdot \bm{\tau} \, \dA.
\label{eq:shear-conjugacy}
\end{equation}

\paragraph{The Romano identification.}

Romano's approach defines the energy-conjugate shear strain implicitly through 
\eqref{eq:shear-conjugacy}. Specifically, $\bm{\gamma}_\perp$ is defined such that:
\begin{equation}
\bm{\gamma}_\perp = \mathbf{B}_\theta\,\bm{b}_{\Section},
\label{eq:romano-shear}
\end{equation}
where $\mathbf{B}_\theta$ is determined by the condition that variations in 
$\bm{b}_{\Section}$ satisfy \eqref{eq:shear-conjugacy}.

From the shear work calculation:
\begin{align*}
\int_\Section \delta\PointShearStrain \cdot \bm{\tau} \, \dA 
&= (\CenterShearLocation \cdot \delta\bm{b}_{\Section})\,T 
+ \delta\bm{b}_{\Section} \cdot \mathbf{H}_W,
\end{align*}
where we used the fact that the $\nabla\bm{\WarpShearIesan}$ term integrates to zero.

However, the torque $T$ under pure flexure loading is:
\begin{equation}
T = G\Jo\,(a_\vartheta + c_\vartheta) = G\Jo\,c_\vartheta \quad \text{(for } a_\vartheta = 0\text{)}.
\end{equation}

This reveals a fundamental coupling: flexure ($\bm{b}_{\Section}$) induces both shear 
force $\bm{F}$ and torque $T$ through the coupling $c_\vartheta = \CenterShearLocation \cdot \bm{b}_{\Section}$.

\paragraph{Complete energy-conjugate trace.}

To achieve full energy conjugacy, we must account for all couplings. 
Define:

The energy-conjugate beam strains are:
\begin{align}
\epsilon &= a_\varepsilon + x\,(-\CentroidLocation \cdot \bm{b}_{\Section}), 
\label{eq:ec-axial} \\
\bm{\gamma}_\perp &= \mathbf{H}^{-\top}\mathbf{H}_R^\top\,\bm{b}_{\Section}, 
\label{eq:ec-shear} \\
\varkappa &= a_\vartheta + c_\vartheta = a_\vartheta + \CenterShearLocation \cdot \bm{b}_{\Section}, 
\label{eq:ec-twist} \\
\bm{\kappa}_\perp &= -\FrameBasis^\times\,\bm{a}_\Omega - x\,(\FrameBasis \times \bm{b}_{\Section}),
\label{eq:ec-curv}
\end{align}
where $\mathbf{H} = \mathbf{H}_\psi + \mathbf{H}_b$ is the total bishear tensor with:
\begin{align}
\mathbf{H}_\psi &:= G \int_\Section (\nabla\bm{\WarpShearIesan})^\top\,(\nabla\bm{\WarpShearIesan}) \, \dA, 
\label{eq:def-H-psi} \\
\mathbf{H}_b &:= G \int_\Section \bm{W}_{(b)}^\top\,(\nabla\bm{\WarpShearIesan}) \, \dA,
\label{eq:def-H-b}
\end{align}
and $\mathbf{H}_R$ is the resultant-shear tensor:
\begin{equation}
\mathbf{H}_R := G \int_\Section (\nabla\bm{\WarpShearIesan}) \, \dA 
+ G \int_\Section \bm{W}_{(b)} \, \dA.
\label{eq:def-H-R}
\end{equation}

\begin{remark}
The expression $\mathbf{H}^{-\top}\mathbf{H}_R^\top$ arises from requiring that 
$\delta\bm{\gamma}_\perp \cdot \bm{F} = \int \delta\PointShearStrain \cdot \bm{\tau}$ 
when $\bm{F}$ and $\bm{\tau}$ are expressed in terms of $\bm{b}_{\Section}$. This is 
Romano's shear correction tensor.
\end{remark}

\paragraph{Assembly of trace matrices.}

From \eqref{eq:ec-axial}--\eqref{eq:ec-curv}, the beam strains 
$\bm{e}(x) = (\epsilon, \bm{\gamma}_\perp, \varkappa, \bm{\kappa}_\perp)^\top$ are:
\begin{equation}
\bm{e}(x) = \TraceDeDpRoot\,\bm{p} + x\,\TraceDeDpOne\,\bm{p},
\end{equation}
with:

\begin{equation}
\TraceDeDpRoot = \begin{pmatrix}
1 & \mathbf{0}^\top & 0 & \mathbf{0}^\top \\[4pt]
\mathbf{0} & \mathbf{0} & \mathbf{0} & \mathbf{H}^{-\top}\mathbf{H}_R^\top \\[4pt]
0 & \mathbf{0}^\top & 1 & \CenterShearLocation^\top \\[4pt]
\mathbf{0} & -\FrameBasis^\times & \mathbf{0} & \mathbf{0}
\end{pmatrix}
\label{eq:T0-romano}
\end{equation}

\begin{equation}
\TraceDeDpOne = \begin{pmatrix}
0 & \mathbf{0}^\top & 0 & -\CentroidLocation^\top \\[4pt]
\mathbf{0} & \mathbf{0} & \mathbf{0} & \mathbf{0} \\[4pt]
0 & \mathbf{0}^\top & 0 & \mathbf{0}^\top \\[4pt]
\mathbf{0} & \mathbf{0} & \mathbf{0} & -\FrameBasis^\times
\end{pmatrix}
\label{eq:T1-romano}
\end{equation}

\paragraph{Warping index.}

The warping index is:
\begin{equation}
w = 6 - \mathrm{rank}\begin{pmatrix} \TraceDeDpRoot \\ \TraceDeDpOne \end{pmatrix}.
\end{equation}

Examining the columns of $\mathbf{T}_{\rm stack}$:
\begin{itemize}
\item Column 1 ($a_\varepsilon$): $(1, \mathbf{0}, 0, \mathbf{0}; 0, \mathbf{0}, 0, \mathbf{0})^\top$ — nonzero.
\item Columns 2--3 ($\bm{a}_\Omega$): $(\mathbf{0}, \mathbf{0}, \mathbf{0}, -\FrameBasis^\times; \ldots)^\top$ — 
the $-\FrameBasis^\times$ block has rank 2, so these contribute 2 independent directions.
\item Column 4 ($a_\vartheta$): $(0, \mathbf{0}, 1, \mathbf{0}; 0, \mathbf{0}, 0, \mathbf{0})^\top$ — nonzero.
\item Columns 5--6 ($\bm{b}_{\Section}$): From $\TraceDeDpRoot$, we get $\mathbf{H}^{-\top}\mathbf{H}_R^\top$ 
in the shear block and $\CenterShearLocation^\top$ in the twist block. From $\TraceDeDpOne$, 
we get $-\CentroidLocation^\top$ in the axial block and $-\FrameBasis^\times$ in the curvature block.
\end{itemize}

For generic sections (where $\mathbf{H}^{-\top}\mathbf{H}_R^\top$ is invertible), 
all six columns are independent.

\begin{proposition}[Warping index of energy-conjugate trace]
\label{prop:warp-index-romano}
For sections with non-degenerate shear response (invertible $\mathbf{H}$), the 
energy-conjugate trace has warping index $w = 0$.
\end{proposition}

\paragraph{Symmetry of the constitutive matrix.}

A key advantage of the energy-conjugate trace is that the resulting section 
stiffness matrix is symmetric. Define the resultant vector 
$\bm{s} = (N, \bm{F}, T, \bm{M})^\top$ and strain vector 
$\bm{e} = (\epsilon, \bm{\gamma}_\perp, \varkappa, \bm{\kappa}_\perp)^\top$. 
The constitutive relation:
\begin{equation}
\bm{s} = \mathbf{C}\,\bm{e}_0
\end{equation}
has $\mathbf{C} = \mathbf{C}^\top$ when the strains are defined by 
\eqref{eq:ec-axial}--\eqref{eq:ec-curv}.

This follows from the variational structure: if $\bm{e}$ and $\bm{s}$ are 
work-conjugate, and the strain energy is a quadratic form in $\bm{e}$, then 
the Hessian (stiffness matrix) is automatically symmetric.

\begin{remark}[Comparison with section-average trace]
The section-average trace and energy-conjugate trace yield different shear 
strain measures:
\begin{align*}
\text{Section-average:} \quad \bar{\bm{\gamma}}_\perp &= \bm{\Gamma}_0\,\bm{b}_{\Section}, \\
\text{Energy-conjugate:} \quad \bm{\gamma}_\perp^{\rm conj} &= \mathbf{H}^{-\top}\mathbf{H}_R^\top\,\bm{b}_{\Section}.
\end{align*}
These are related by a gauge transformation. The energy-conjugate shear 
resultant $\bm{F}^{\rm conj} = \bm{F}$ (the classical shear force), while the 
section-average shear resultant $\bar{\bm{F}}$ differs from $\bm{F}$ by terms 
involving the bishear tensors.
\end{remark}

\begin{remark}[Physical interpretation]
The energy-conjugate trace does not correspond to any simple geometric 
operation on the displacement field (unlike the pointwise or section-average 
traces). The ``rotation'' $\bm{\theta}$ implicit in this trace is defined 
indirectly through the requirement that $\bm{\gamma} = \bm{u}' + \FrameBasis \times \bm{\theta}$ 
satisfy energy conjugacy.

Consequently, the energy-conjugate trace may not satisfy the plane-section 
property: for a section-rigid displacement, the extracted ``rotation'' may 
include contributions that would conventionally be attributed to translation 
or warping.
\end{remark}

\begin{remark}[Boundary conditions under energy-conjugate trace]
Care must be taken when imposing boundary conditions with the energy-conjugate 
trace. 
A ``clamped'' condition $\bm{u}(0) = \bm{\theta}(0) = \mathbf{0}$ does not 
correspond to fixing the physical displacement and rotation of any particular 
point or average, but rather to a condition on the energy-conjugate kinematic 
measures. 
This distinction becomes significant when comparing beam solutions 
to three-dimensional analyses.
\end{remark}

